\documentclass{article}
%\usepackage{assymbb}
\begin{document}
Espace a 2 dim borne avec $x \in [0, 1]$ et $y \in [0, 1] $\\

Soient $A(a_x, a_y)$, $B(b_y, b_y)$ deux points distincts pour lesquels on souhaite connaitre la médiatrice med\\

Soit $C(c_x, c_y)$ le point dont on souhaite calculer la distance à la médiatrice\\

Soit $M(m_x, m_y)$ le milieu du segment [AB]. Par def $M \in med$\\
$m_x = (a_x+b_x)/2$\\
$m_y = (a_y+b_y)/2$\\

Soit $\vec{AB}$ le vecteur transformant A vers B\\
$\vec{AB}_x = b_x-a_x$\\
$\vec{AB}_y = b_y-a_y$\\

Soit $\vec{p}$ un vecteur perpendiculaire à $\vec{AB}$\\
$\vec{p}_x = -\vec{AB}_y = a_y-b_y$\\
$\vec{p}_y = \vec{AB}_y = b_x-a_x$\\

Par définition, ce vecteur est colinéaire à med.\\
On peut définir med $(M, \vec{p})$ car perpendiculaire au segment [AB] et coupant par son milieu\\

Soit $H(h_x, h_y)$ la projection orthogonale de C sur med\\
H doit satisfaire 2 conditions:\\
(1) $\vec{CH} \cdot \vec{p} = 0$ (perpendiculaire)\\
(2) $H \in (M,\vec{p})$\\

On précise:\\
$\vec{CH}_x = h_x-c_x$\\
$\vec{CH}_y = h_y-c_y$\\

Developpons (2):\\
$\exists t \in R$ tq $H = M + t*\vec{p}$\\
Donc\\
$h_x = m_x + t*\vec{p}_x = (a_x+b_x)/2 + t*(a_y - b_y)$ (4)\\
$h_y = m_y + t*\vec{p}_y = (a_y+b_y)/2 + t*(b_x - a_x)$ (5)\\

D'apres (1):\\
$(h_x-c_x)(a_y-b_y) + (h_y-c_y)(b_x-a_x) = 0$\\
$\Leftrightarrow h_x(a_y-b_y) - c_x(a_y - b_y) + h_y(b_x - a_x) - c_y(b_x - a_x) = 0$\\
$\Leftrightarrow (\frac{a_x+b_x}{2} + t(a_y - b_y))(a_y-b_y) - c_x(a_y - b_y) + (\frac{a_y+b_y}{2} + t(b_x - a_x))(b_x - a_x) - c_y(b_x - a_x) = 0$\\
...\\
$\Leftrightarrow t = \frac{c_x(a_y-b_y)+c_y(b_x-a_x)+a_yb_y-a_yb_x}{(a_y-b_y)^2 + (b_x-a_x)^2}$\\

\end{document}